\chapter*{摘要}
\addcontentsline{toc}{chapter}{摘要}

本项目面向废塑料瓶分拣场景，使用 ALOHA 双臂机器人学习一个完整长流程：左侧工作臂从散乱瓶堆中取瓶并稳定瓶身，右侧工作臂抓住瓶盖并旋下；随后左臂将瓶身投入一个回收盒，右臂将瓶盖投入另一个回收盒。任务同时包含随机目标发现、双臂协作、精细接触、连续旋转、物体分离和分类投放，比“一次抓起一个物体”明显更复杂。

第一年度已经形成从专用采集软件、数据编辑中心、数据版本发布、模型训练到机器人连续运行的完整工程闭环。数据平台现有 51 个有效 ALOHA 项目、2,413 条轨迹、2,907,804 帧，按记录帧率换算约 16.15 小时；正式部署模型实际使用其中 25 个唯一数据仓库、1,051 条轨迹、879,852 帧，约 4.89 小时，过滤后可训练部分为 844,102 帧、约 4.69 小时。\ev{E03,E04} 全量数值检查未发现非有限状态/动作、全零动作、轨迹内时间倒退或帧号中断；视频尚未逐帧全部解码，因此不能称为“无任何数据问题”。\ev{E05}

正式模型接收顶部总览、左腕和右腕三路画面，以及双臂和夹爪的 14 个状态量；每次产生未来 50 个控制时刻的双臂动作，运行控制频率为 50 Hz，并在当前动作片段执行期间提前准备下一段动作。\ev{E06,E07} 保存模型包含 838,358,468 个参数。训练历史有 6,059 条记录，覆盖 0--59,990 步，跨度约 51.4 小时；训练误差从 0.0677 降至 0.000671。现场实际加载的是第 19,000 步保存版本，而不是最终训练记录。\ev{E08,E09} 训练记录没有验证集、测试集或机器人任务成功率，因而误差下降只能说明模型更接近示教动作，不能代替真机成功率。

现场展示中观察到机器人完成完整分拣循环，并有连续工作超过 1 小时、平均约 2 瓶/分钟的表现；同时观察到一定的桌面位置和瓶子条件变化适应能力。\ev{E10} 但是没有完整录像、逐瓶计数表、固定测试协议或自动成功判定，因此上述时长和吞吐量均为\textbf{现场观测估计}，不能写成正式基准或稳定成功率。当前最清楚的失败是：无盖瓶仍会执行空拧，部分瓶子会被倒着抓起。进一步审计发现，6 个明确为无盖条件的数据仓库、158 条轨迹仍全部使用无条件“拧盖”指令；这与空拧现象一致，是高优先级数据风险，但尚不能证明唯一因果关系。\ev{E11}

团队还开展了强化学习方向研究：835 条冲洗任务轨迹、238,816 个训练片段、连续 5 个研究轮次以及 6,154 万参数的动作优化模型已经可复核；离线验证动作误差相对下降约 6.7\%。\ev{E15,E16} 由于缺少同条件基础模型对照和真机测试，该研究目前只能称为“训练闭环已建立”，不能称为“已经提升真机产能”。

下一年度的顺序应为：第一，建立固定场景、逐瓶计数、阶段判定和失败分类，先让结果可信；第二，补采“无盖时跳过拧盖”、正反方向、抓取纠错和接触阶段数据，提升基础模仿学习模型；第三，在同一评测协议下验证强化学习增益；第四，推进把空瓶插入水管的数字孪生、视觉定位与接触控制研究。所谓毫米级精度目前是任务设计目标，尚需用真实几何标定和重复插入误差来验收。

\figfull[.97\textwidth]{evidence_grade.pdf}{成果证据分级。可复核事实、现场观测估计和下一阶段量化指标采用不同等级呈现，训练误差或单次展示不作为稳定成功率。}

\section*{一页成果概览}
\begin{table}[H]
\centering
\caption{第一年度可确认成果与边界}
\begin{tabularx}{\textwidth}{>{\bfseries}p{31mm}p{35mm}X}
\toprule
方面 & 已取得成果 & 当前边界\\
\midrule
完整任务 & 观察到取瓶、旋盖、瓶/盖分箱的连续循环 & 无逐瓶计数和完整录像，不能给正式成功率\\
连续运行 & 现场估计超过 1 小时、约 2 瓶/分钟 & 不是带置信区间的标准吞吐量\\
数据资产 & 51 个有效项目、2,413 条轨迹、约 16.15 小时 & 正式模型只使用约 4.89 小时唯一子集\\
正式模型 & 8.38 亿参数、三路视觉、双臂连续动作 & 训练误差无验证/测试对照\\
工程迭代 & 41 次可追溯运行尝试、14 组配置 & 大量运行中断，不能称为 41 个成功实验\\
可解释记录 & 8,223 份真实注意力记录 & 无结果标签与遮挡实验，不能作因果解释\\
强化学习 & 数据、保存模型和离线验证闭环已形成 & 未证明真机优于基础模型\\
\bottomrule
\end{tabularx}
\end{table}
